\cotefile{157-5}
\scanpage{618}
\mspage{6}
alors a) on a $\Phi_1 \mathrel{\longleftarrow\!\mid\!\longrightarrow} \Psi$ et b) $\Phi_1^{\sim}=\Phi^{\sim}=\Phi_1^{\natural}$.

Ce corollaire résulte du théorème en vertu de
\uline{$\overline{\mathrm{XV}}$} prop. 4.2, dont je rappelle ici l'essentiel

\uline{Proposition} 1.3. Soit $\Phi \mathrel{\longleftarrow\!\mid\!\longrightarrow} \Psi$. Alors
$\Phi^{\sim}=\Phi^{\natural}=\Psi^{*}$, et dualement on a $\overline{\Psi}=\Psi^{\natural}=\Phi_{\unsure{*|\star}}$.

Notons aussi, \unsure{dans|en} cet ordre d'idées

\uline{\deleted{Proposition}\inserted{Corollaire} 1.4} Soit $\Phi\longleftrightarrow\Psi$, et soit $W$
une partie \deleted{\unsure{quelconque}}\inserted{quelconque} de $\mathrm{Fl}\,M$. \deleted{\unsure{et W, pour}}

Pour qu'on ait $\Phi\subset W^{\natural}$, il suffit que
toute flèche de $M$ se factorise en $f=pi$,
avec $i\in W$, $p\in\Psi$ \deleted{\unsure{si}}. (Et cette condition
est \inserted{évidemment} aussi nécessaire si $\Phi \mathrel{\longleftarrow\!\mid\!\longrightarrow} \Psi$, et $W$
stable par facteurs directs.)

On dit qu'une partie $\Phi_0$ de $\Phi$ est \deleted{une}
$Q$-\uline{génératrice}\inserted{\unsure{à g.}} de $\Phi$ si elle satisfait les conditions

\deleted{\unsure{a)} $\unsure{\Phi_0}$ \unsure{est} \illegible{1}}

\[
\deleted{\unsure{b)}}\ \Phi\subset\Phi_0^{\sim},\tag*{(1.9)}
\]

et on dit que $\Phi$ est une partie \uline{spéciale} \deleted{\illegible{1}}
si $\Phi$ admet une \inserted{\uuline{petite}} partie $Q$-génératrice à gauche.
Si $\Phi_0\subset\Phi$ est petit, et si $M$ satisfait les
conditions du th. 1.1, alors $\Phi_0$, est une
partie $Q$-génératrice (à gauche) sss son
$Q$-saturé total (à g.) $\Phi_0'''$ contient $\Phi$
(condition en apparence plus forte que $\Phi\subset\Phi_0^{\sim}$,
puisque en général on a $\Phi_0'''\subset\Phi_0^{\sim}$). Cela
signifie donc \unsure{aussi|que}, via 1.2, que \deleted{\unsure{toute}}
% archive sheet number: 617
